% ============================================================
% 测试卷 — 第1章 集合与常用逻辑用语
% ============================================================
\documentclass[12pt, a4paper]{ctexart}
\usepackage[margin=2cm]{geometry}
\usepackage{amsmath, amssymb}
\usepackage{enumitem}
\usepackage{fancyhdr}
\usepackage{xcolor}
\usepackage{tikz}
\usepackage{graphicx}
\usepackage{tabularx}
\usepackage{multicol}

\definecolor{mainblue}{RGB}{41, 98, 168}

\pagestyle{fancy}
\fancyhf{}
\fancyhead[L]{\small\textcolor{mainblue}{2025-2026学年度·第1章测试卷}}
\fancyhead[R]{\small\textcolor{mainblue}{数学（必修第一册）}}
\fancyfoot[C]{\small —\ \thepage\ —}
\renewcommand{\headrulewidth}{0.6pt}
\renewcommand{\headrule}{\hbox to\headwidth{\color{mainblue}\leaders\hrule height \headrulewidth\hfill}}

\setCJKmainfont{Hiragino Sans GB}
\setCJKsansfont{Heiti SC}

\newcommand{\blank}[1]{\underline{\hspace{#1}}}

\begin{document}

\begin{center}
{\Large\bfseries 第1章\quad 集合与常用逻辑用语\quad 测试卷}\\[6pt]
{\normalsize 数学（必修第一册）}\\[4pt]
{\small 考试时间：120分钟\quad 满分：150分}
\end{center}

\vspace{2pt}
\begin{center}
\begin{tabular}{|c|c|c|c|c|}
\hline
题号 & 一 & 二 & 三 & 总分 \\ \hline
得分 & & & & \\ \hline
\end{tabular}
\end{center}
\vspace{4pt}

\noindent\textbf{注意事项：}
\begin{enumerate}[leftmargin=2em]
  \item 答题前，请在答题卡上填写姓名、班级和考号。
  \item 选择题请用2B铅笔在答题卡上涂黑对应选项。
  \item 非选择题请用黑色签字笔在答题卡指定区域作答。
\end{enumerate}

\vspace{6pt}

% ============ 一、单项选择题 ============
\noindent\rule{\textwidth}{0.4pt}\\[2pt]
{\large\bfseries 一、单项选择题}（本大题共8小题，每小题5分，共40分。每小题只有一个选项符合题目要求。）

\vspace{6pt}

\begin{enumerate}
\item 设集合 $A=\{x\mid x^2-3x+2=0\}$，$B=\{1,2,3\}$，则 $A\cap B=$（\quad）
  \begin{enumerate}[label=\Alph*., itemsep=0pt]
    \item $\{1\}$ \quad \item $\{2\}$ \quad \item $\{1,2\}$ \quad \item $\{1,2,3\}$
  \end{enumerate}

\item 已知全集 $U=\{1,2,3,4,5\}$，集合 $A=\{1,3,5\}$，则 $\complement_U A=$（\quad）
  \begin{enumerate}[label=\Alph*., itemsep=0pt]
    \item $\{1,3,5\}$ \quad \item $\{2,4\}$ \quad \item $\{1,2\}$ \quad \item $\{4,5\}$
  \end{enumerate}

\item 设 $p: x>1$，$q: x>2$，则 $p$ 是 $q$ 的（\quad）
  \begin{enumerate}[label=\Alph*., itemsep=0pt]
    \item 充分不必要条件 \quad \item 必要不充分条件
    \item 充要条件 \quad \item 既不充分也不必要条件
  \end{enumerate}

\item 已知集合 $A=\{x\mid -1<x<3\}$，$B=\{x\mid x>a\}$，若 $A\subseteq B$，则 $a$ 的取值范围是（\quad）
  \begin{enumerate}[label=\Alph*., itemsep=0pt]
    \item $a\leq -1$ \quad \item $a<-1$ \quad \item $a\leq 3$ \quad \item $a<3$
  \end{enumerate}

\item 命题"$\forall x\in\mathbb{R}$，$x^2+1>0$"的否定是（\quad）
  \begin{enumerate}[label=\Alph*., itemsep=0pt]
    \item $\forall x\in\mathbb{R}$，$x^2+1\leq 0$
    \item $\exists x_0\in\mathbb{R}$，$x_0^2+1>0$
    \item $\exists x_0\in\mathbb{R}$，$x_0^2+1\leq 0$
    \item $\forall x\in\mathbb{R}$，$x^2+1<0$
  \end{enumerate}

\item 设集合 $A=\{1,a,a^2\}$，若 $2\in A$，则 $a=$（\quad）
  \begin{enumerate}[label=\Alph*., itemsep=0pt]
    \item $2$ \quad \item $-1$ \quad \item $2$ 或 $-1$ \quad \item $1$
  \end{enumerate}

\item 已知 $p: |x-1|<2$，$q: x(x-3)<0$，则 $p$ 是 $q$ 的（\quad）
  \begin{enumerate}[label=\Alph*., itemsep=0pt]
    \item 充分不必要条件 \quad \item 必要不充分条件
    \item 充要条件 \quad \item 既不充分也不必要条件
  \end{enumerate}

\item 设集合 $A=\{x\mid x^2-5x+6=0\}$，$B=\{x\mid mx-1=0\}$，若 $B\subseteq A$，则 $m$ 的所有可能值组成的集合为（\quad）
  \begin{enumerate}[label=\Alph*., itemsep=0pt]
    \item $\left\{0,\dfrac{1}{2},\dfrac{1}{3}\right\}$ \quad \item $\left\{\dfrac{1}{2},\dfrac{1}{3}\right\}$
    \item $\left\{0,\dfrac{1}{2}\right\}$ \quad \item $\left\{\dfrac{1}{2}\right\}$
  \end{enumerate}
\end{enumerate}

\vspace{6pt}

% ============ 二、多项选择题 ============
\noindent\rule{\textwidth}{0.4pt}\\[2pt]
{\large\bfseries 二、多项选择题}（本大题共4小题，每小题5分，共20分。全部选对得5分，部分选对得2分，有选错的得0分。）

\vspace{6pt}

\begin{enumerate}[resume]
\item 下列命题中正确的是（\quad）
  \begin{enumerate}[label=\Alph*., itemsep=0pt]
    \item 空集是任何集合的子集
    \item $\{0\}$ 是空集
    \item 若 $A\cup B=A$，则 $B\subseteq A$
    \item 若 $A\cap B=\varnothing$，则 $A=\varnothing$ 或 $B=\varnothing$
  \end{enumerate}

\item 已知集合 $A=\{x\mid x^2-2x-3<0\}$，$B=\{x\mid x>1\}$，则（\quad）
  \begin{enumerate}[label=\Alph*., itemsep=0pt]
    \item $A=(-1,3)$ \quad \item $A\cap B=(1,3)$
    \item $A\cup B=(-1,+\infty)$ \quad \item $\complement_{\mathbb{R}}A=(-\infty,-1]\cup[3,+\infty)$
  \end{enumerate}

\item 下列说法正确的是（\quad）
  \begin{enumerate}[label=\Alph*., itemsep=0pt]
    \item "$x=1$"是"$x^2-1=0$"的充分条件
    \item "$x>1$"是"$x^2>1$"的充分条件
    \item "$x=y$"是"$\sin x=\sin y$"的必要条件
    \item "$ab=0$"是"$a=0$"的必要条件
  \end{enumerate}

\item 设全集 $U=\mathbb{R}$，集合 $A=\{x\mid |x|<2\}$，$B=\{x\mid x^2-4x+3\leq 0\}$，则（\quad）
  \begin{enumerate}[label=\Alph*., itemsep=0pt]
    \item $A=(-2,2)$ \quad \item $B=[1,3]$
    \item $A\cap B=[1,2)$ \quad \item $A\cup B=(-2,3]$
  \end{enumerate}
\end{enumerate}

\vspace{6pt}

% ============ 三、填空题 ============
\noindent\rule{\textwidth}{0.4pt}\\[2pt]
{\large\bfseries 三、填空题}（本大题共4小题，每小题5分，共20分。）

\vspace{6pt}

\begin{enumerate}[resume]
\item 设集合 $A=\{1,2,3\}$，$B=\{2,3,4,5\}$，则 $A\cup B=$ \underline{\hspace{3cm}}。

\item 已知集合 $A=\{(x,y)\mid x+y=3\}$，$B=\{(x,y)\mid x-y=1\}$，则 $A\cap B=$ \underline{\hspace{3cm}}。

\item 命题"$\exists x_0\in\mathbb{R}$，$x_0^2+2x_0+2\leq 0$"的真假为\underline{\hspace{2cm}}（填"真命题"或"假命题"）。

\item 设集合 $A=\{x\mid x^2-3x+2=0\}$，$B=\{x\mid x^2-ax+a-1=0\}$，若 $B\subseteq A$，则 $a$ 的取值范围为 \underline{\hspace{3cm}}。
\end{enumerate}

\vspace{6pt}

% ============ 四、解答题 ============
\noindent\rule{\textwidth}{0.4pt}\\[2pt]
{\large\bfseries 四、解答题}（本大题共6小题，共70分。解答应写出文字说明、证明过程或演算步骤。）

\vspace{6pt}

\begin{enumerate}[resume]
\item[17.]（本小题10分）\\
已知集合 $A=\{x\mid x^2-5x+6=0\}$，$B=\{x\mid mx+1=0\}$。
\begin{enumerate}[(1)]
  \item 求集合 $A$；
  \item 若 $A\cup B=A$，求实数 $m$ 的值。
\end{enumerate}

\vspace{6cm}

\item[18.]（本小题12分）\\
已知集合 $A=\{x\mid 2\leq x<5\}$，$B=\{x\mid 3<x<7\}$。
\begin{enumerate}[(1)]
  \item 求 $A\cap B$，$A\cup B$；
  \item 求 $\complement_{\mathbb{R}}(A\cap B)$。
\end{enumerate}

\vspace{6cm}

\item[19.]（本小题12分）\\
判断下列命题的真假，并说明理由：
\begin{enumerate}[(1)]
  \item $\forall x\in\mathbb{R}$，$x^2-2x+3>0$；
  \item $\exists x_0\in\mathbb{Q}$，$x_0^2=2$；
  \item 若 $a>b$，则 $a^2>b^2$。
\end{enumerate}

\vspace{6cm}

\item[20.]（本小题12分）\\
已知 $p: x^2-8x-20\leq 0$，$q: x^2-2x+1-m^2\leq 0$（$m>0$）。若 $p$ 是 $q$ 的充分不必要条件，求实数 $m$ 的取值范围。

\vspace{7cm}

\item[21.]（本小题12分）\\
设集合 $A=\{x\mid x^2-3x+2=0\}$，$B=\{x\mid x^2+2(a+1)x+a^2-5=0\}$。
\begin{enumerate}[(1)]
  \item 若 $A\cap B=\{2\}$，求实数 $a$ 的值；
  \item 若 $A\cup B=A$，求实数 $a$ 的取值范围。
\end{enumerate}

\vspace{7cm}

\item[22.]（本小题12分）\\
已知集合 $A=\{x\mid x^2-ax+a^2-19=0\}$，$B=\{x\mid x^2-5x+6=0\}$，$C=\{x\mid x^2+2x-8=0\}$。
\begin{enumerate}[(1)]
  \item 若 $A\cap B=A\cup B$，求 $a$ 的值；
  \item 若 $A\cap B\neq\varnothing$，$A\cap C=\varnothing$，求 $a$ 的值。
\end{enumerate}

\vspace{7cm}
\end{enumerate}

\vfill
\begin{center}
\small\textcolor{gray}{本试卷仅供教学参考}
\end{center}

\end{document}
